%% =====================================================================
%%  T-09-04  The Envy Margin
%%  -------------------------------------------------------------------
%%  One idea per file. Core is mandatory. Slots in canonical order.
%%  Max 700 words. No layout commands. No filler. New source material
%%  for this entry goes in sources/B3/T-09-04.tex -- never by
%%  rewriting this file (docs/RULES.md R0.1).
%%  =====================================================================
%% @kind: topic
%% @id: T-09-04
%% @title: The Envy Margin
%% @subtitle: Never appear too perfect --- envy recruits silent enemies
%% @chapter: c09
%% @order: 40
%% @status: stable
%% @sources: B3
%% @updated: 2026-09-19

\topic{T-09-04}{The Envy Margin}{Never appear too perfect --- envy recruits silent enemies}

\begin{core}
Reputation work has an upper bound: appearing better than others is dangerous, and appearing to have no faults at all is the most dangerous position there is. Envy does not argue --- it recruits silent enemies who wait for your first fall. The discipline is to keep a margin of visible imperfection: occasional defects shown, harmless vices admitted, so admiration arrives without its tax.
\end{core}
\begin{mechanism}
B3 states the rule plainly: the most dangerous of all appearances is having no weaknesses, because envy creates silent enemies --- people who never announce the hostility but celebrate your downfall and quietly work toward it. The mechanism is comparison pain: unblemished success reads, to people who cannot name the feeling, as an accusation of their own lives; the resentment has no legitimate outlet, so it goes underground and waits. Displayed defects discharge it --- a named flaw, a self-deprecating story, an admitted vice that costs nothing strategic converts the perfect stranger into a flawed human at human distance, and comparison pain becomes affinity. The source's epigram closes the case: perfection is survivable only for gods and for the dead. The margin is also armour: an enemy armed with your confessed flaw is carrying stolen ammunition, and the reveal's sting is already spent.
\end{mechanism}
\begin{conditions}
Strongest where audiences compare themselves to you directly --- same office, same profession, same circle --- and rises with your proximity to their level. Weakens for distant figures (audiences envy the near, not the legendary) and in formal evaluation settings where confessed defects are priced literally. The margin must be genuine: performed humility is detected and read as a deeper form of superiority.
\end{conditions}
\begin{application}
  \item Audit the image you present to those nearest your position. If it is flawless, it is a target.
  \item Keep two or three real, harmless defects on display --- a vice admitted, a skill lacked, a story where you were the fool.
  \item Never fake the defects. A manufactured weakness reads as condescension and costs more than perfection.
  \item When praised above your peers, redistribute: credit the team, name the luck, in plain words, without protest.
  \item Watch your own envy. The resentments you hide are the working model of the ones aimed at you.
\end{application}
\begin{feedback}
  \item Working: peers tease you and include you. Teasing is the tax receipt --- the margin is holding.
  \item Working: your wins draw congratulations instead of silence. Silent congratulations are the flag.
  \item Failing: a confessed defect cost you standing with evaluators. Wrong audience --- the margin is for peers, not examiners.
  \item Failing: people below you have stopped bringing you their problems. You became too perfect to approach.
\end{feedback}
\begin{failure}
The margin overdrawn becomes the reputation: defects displayed past a point erode the assets the reputation chapters built (T-09-01). Timing has its own trap --- vices confessed while competitors circle hand them the weapon pre-wrapped. And the rule curdles into deliberate underachievement, which is not envy management but fear wearing modesty's coat.
\end{failure}
\begin{countermeasures}
  \item Notice your own silent envy: whoever's success you cannot congratulate warmly is teaching you this mechanism from inside.
  \item Distrust the flawless. A colleague with no visible weaknesses is hiding them, or hiding the work (T-13-03).
  \item When someone confesses a harmless vice to you, ask which comparison the confession is managing. The deflection has a target --- sometimes you.
  \item In evaluation, price results, not polish. The perfect interview is a costume with a tailor (T-23-03).
\end{countermeasures}

\sources{B3}
\seesources
\seealso{T-09-01, T-09-02, T-12-04}
